\documentclass{article}
\usepackage[utf8]{inputenc}
\usepackage[german]{babel}
\usepackage{amsmath, amssymb, amsthm}
\usepackage{booktabs}
\usepackage{array}

\newtheorem{theorem}{Satz}
\newcolumntype{L}[1]{>{\raggedright\arraybackslash}p{#1}}

\title{Diskussion der Quantoren-Verschmelzung}
\author{Nach Daniel J. Velleman: \emph{How to Prove It}}
\date{\today}

\begin{document}

\maketitle

\section*{Das mathematische Problem}

Wir wollen beweisen, dass die verschachtelte Aussage (vor der Verschmelzung) logisch äquivalent zur einfachen Aussage (nach der Verschmelzung) ist. Das heißt, für feste Mengen $x$ und $\mathcal{F}$ gilt:
\[
\exists Y \, \Big( \exists A \, (A \in \mathcal{F} \land \mathcal{P}(A) = Y) \land x \in Y \Big)
\iff
\exists A \, (A \in \mathcal{F} \land x \in \mathcal{P}(A))
\]

Da es sich um eine Äquivalenz ($\iff$) handelt, teilt sich unser Velleman-Scratch-Work in zwei Richtungen auf: Hinrichtung ($\rightarrow$) und Rückrichtung ($\leftarrow$).

---

\section*{Richtung 1: Hinrichtung ($\rightarrow$)}
Wir nehmen die verschachtelte Aussage als bekannt an und wollen die verschmolzene Aussage zeigen.

\subsection*{Scratch Work (Hinrichtung)}
\begin{center}
\small
\begin{tabular}{ L{0.45\textwidth} | L{0.45\textwidth} }
\toprule
\textbf{Givens} & \textbf{Goals} \\
\midrule
% Schritt 0
1. $\exists Y \, \Big( \exists A \, (A \in \mathcal{F} \land \mathcal{P}(A) = Y) \land x \in Y \Big)$ &
\emph{Goal:} $\exists A \, (A \in \mathcal{F} \land x \in \mathcal{P}(A))$ \\[1ex]
\hline
% Schritt 1: Existenzquantor für Y in den Givens abbauen
2. $Y_0$ ist ein konkretes Objekt. \newline
3. $\exists A \, (A \in \mathcal{F} \land \mathcal{P}(A) = Y_0) \land x \in Y_0$ \newline
(\emph{Rezept: $\exists$ im Given genutzt}) &
\emph{Goal:} $\exists A \, (A \in \mathcal{F} \land x \in \mathcal{P}(A))$ \\[1ex]
\hline
% Schritt 2: UND-Verknüpfung in Given 3 aufspalten
2. $Y_0$ ist ein konkretes Objekt. \newline
4. $\exists A \, (A \in \mathcal{F} \land \mathcal{P}(A) = Y_0)$ \newline
5. $x \in Y_0$ \newline
(\emph{Rezept: $\land$ im Given genutzt}) &
\emph{Goal:} $\exists A \, (A \in \mathcal{F} \land x \in \mathcal{P}(A))$ \\[1ex]
\hline
% Schritt 3: Existenzquantor für A in Given 4 abbauen
2. $Y_0$ ist ein konkretes Objekt. \newline
5. $x \in Y_0$ \newline
6. $A_0$ ist ein konkretes Objekt. \newline
7. $A_0 \in \mathcal{F} \land \mathcal{P}(A_0) = Y_0$ \newline
(\emph{Rezept: $\exists$ im Given genutzt}) &
\emph{Goal:} $\exists A \, (A \in \mathcal{F} \land x \in \mathcal{P}(A))$ \\[1ex]
\hline
% Schritt 4: UND-Verknüpfung in Given 7 aufspalten und substituieren
2. $Y_0$ ist ein konkretes Objekt. \newline
5. $x \in Y_0$ \newline
6. $A_0$ ist ein konkretes Objekt. \newline
8. $A_0 \in \mathcal{F}$ \newline
9. $\mathcal{P}(A_0) = Y_0$ \newline
10. $x \in \mathcal{P}(A_0)$ \quad (Substituiere 9 in 5) &
\emph{Goal:} $\exists A \, (A \in \mathcal{F} \land x \in \mathcal{P}(A))$ \\[1ex]
\hline
% Schritt 5: Existenzquantor im Goal erfüllen (Konstruktion!)
% Wir wählen A_0 als unseren Zeugen für das Goal.
[Givens wie oben in Schritt 4] &
\emph{Neues Goal:} $A_0 \in \mathcal{F} \land x \in \mathcal{P}(A_0)$ \newline
(\emph{Rezept: $\exists$ im Goal mit Kandidat $A_0$}) \newline
\newline
\checkmark \quad (Ziel erreicht, da Given 8 und Given 10 genau das fordern!) \\
\bottomrule
\end{tabular}
\end{center}

---

\section*{Richtung 2: Rückrichtung ($\leftarrow$)}
Wir nehmen die einfache, verschmolzene Aussage an und wollen die verschachtelte Struktur rekonstruieren.

\subsection*{Scratch Work (Rückrichtung)}
\begin{center}
\small
\begin{tabular}{ L{0.45\textwidth} | L{0.45\textwidth} }
\toprule
\textbf{Givens} & \textbf{Goals} \\
\midrule
% Schritt 0
1. $\exists A \, (A \in \mathcal{F} \land x \in \mathcal{P}(A))$ &
\emph{Goal:} $\exists Y \, \Big( \exists A \, (A \in \mathcal{F} \land \mathcal{P}(A) = Y) \land x \in Y \Big)$ \\[1ex]
\hline
% Schritt 1: Existenzquantor im Given abbauen
2. $A_0$ ist ein konkretes Objekt. \newline
3. $A_0 \in \mathcal{F} \land x \in \mathcal{P}(A_0)$ &
[Goal wie oben] \\[1ex]
\hline
% Schritt 2: UND-Verknüpfung aufspalten
2. $A_0$ ist ein konkretes Objekt. \newline
4. $A_0 \in \mathcal{F}$ \newline
5. $x \in \mathcal{P}(A_0)$ &
[Goal wie oben] \\[1ex]
\hline
% Schritt 3: Existenzquantor im Goal angreifen (Konstruktion)
% Welches Y löst das Goal? Wir WÄHLEN Y = \mathcal{P}(A_0)!
2. $A_0$ ist ein konkretes Objekt. \newline
4. $A_0 \in \mathcal{F}$ \newline
5. $x \in \mathcal{P}(A_0)$ &
\emph{Neues Goal:} \newline
$\exists A \, (A \in \mathcal{F} \land \mathcal{P}(A) = \mathcal{P}(A_0)) \land x \in \mathcal{P}(A_0)$ \newline
(\emph{Rezept: $\exists$ im Goal mit Kandidat $Y = \mathcal{P}(A_0)$}) \\[1ex]
\hline
% Schritt 4: Das UND-Goal aufteilen (Zweig 1 und Zweig 2)
% Zweig 2 (x \in \mathcal{P}(A_0)) ist durch Given 5 sofort wahr.
% Bleibt nur noch das linke Goal übrig:
2. $A_0$ ist ein konkretes Objekt. \newline
4. $A_0 \in \mathcal{F}$ \newline
5. $x \in \mathcal{P}(A_0)$ &
\emph{Verbleibendes Goal:} \newline
\phantom{X} $\exists A \, (A \in \mathcal{F} \land \mathcal{P}(A) = \mathcal{P}(A_0))$ \\[1ex]
\hline
% Schritt 5: Letzten Existenzquantor im Goal lösen
% Welches A erfüllt dieses Goal? Natürlich unser A_0!
[Givens wie oben] &
\emph{Letztes Goal:} \newline
$A_0 \in \mathcal{F} \land \mathcal{P}(A_0) = \mathcal{P}(A_0)$ \newline
(\emph{Kandidat $A = A_0$}) \newline
\newline
\checkmark \quad (Ziel erreicht! $A_0 \in \mathcal{F}$ ist Given 4, und $\mathcal{P}(A_0) = \mathcal{P}(A_0)$ ist trivialerweise wahr!) \\
\bottomrule
\end{tabular}
\end{center}

\end{document}
